\documentclass[11pt,a4paper]{article}
\usepackage[utf8]{inputenc}
\usepackage[vietnamese]{babel}
\usepackage{amsmath,amssymb}
\usepackage{geometry}
\geometry{a4paper, total={180mm,260mm}, left=15mm, right=15mm, top=18mm, bottom=18mm}
\usepackage{tikz}
\usetikzlibrary{shapes.geometric, arrows.meta, positioning, calc}
\usepackage{tcolorbox}
\tcbuselibrary{skins,breakable}
\usepackage{xcolor}
\usepackage{multicol}
\usepackage{array}
\usepackage{fontawesome}
\usepackage{fancyhdr}

\pagestyle{fancy}
\fancyhf{}
\renewcommand{\headrulewidth}{0.5pt}
\renewcommand{\footrulewidth}{0.5pt}
\fancyhead[L]{\small\bfseries\color{teal}\faBook\ SỔ TAY KIẾN THỨC TOÁN 3}
\fancyhead[R]{\small\bfseries\color{teal}Lớp Toán Thầy Lĩnh -- loptoanthaylinh.vn}
\fancyfoot[C]{\small\color{gray}-- Sổ tay kiến thức: Bài 39. So sánh số lớn gấp mấy lần số bé --}

\definecolor{primary}{HTML}{005F73}
\definecolor{secondary}{HTML}{0A9396}
\definecolor{accent}{HTML}{EE9B00}
\definecolor{danger}{HTML}{AE2012}
\definecolor{boxbg}{HTML}{F4F6F8}

\begin{document}

% --- BANNER TIÊU ĐỀ ---
\begin{tcolorbox}[
  enhanced,
  colback=white,
  colframe=primary,
  arc=3mm,
  boxrule=1.2pt,
  left=4mm, right=4mm, top=2.5mm, bottom=2.5mm,
  drop shadow=black!10!white
]
  \centering
  {\Large\bfseries\color{primary}\faBookmark\ BÀI 39: SO SÁNH SỐ LỚN GẤP MẤY LẦN SỐ BÉ}\\[1.5mm]
  {\small\itshape\color{black!75}Quy tắc lấy số lớn chia cho số bé và liên hệ số bé bằng một phần mấy số lớn trên 1 trang A4}
\end{tcolorbox}

\vspace{-2mm}


\noindent
\begin{minipage}[t]{0.49\textwidth}\vspace{0pt}
\begin{tcolorbox}[
  enhanced,
  colback=blue!4!white,
  colframe=blue!75!black,
  arc=2mm, boxrule=0.8pt,
  left=2mm, right=2mm, top=1.8mm, bottom=1.8mm,
  title={\footnotesize\bfseries\color{white}\faKey\ 1. Quy tắc so sánh gấp số lần}
]
{\footnotesize
\begin{itemize}\setlength{\itemsep}{1.5pt}\setlength{\leftskip}{-2.5mm}
  \item Muốn biết \textbf{số lớn gấp mấy lần số bé}, ta lấy số lớn \textbf{chia cho số bé}:
  $$\text{Số lần} = \text{Số lớn} : \text{Số bé}$$
  \item \textbf{Ví dụ cốt lõi:}
  \begin{itemize}\setlength{\itemsep}{1pt}\setlength{\leftskip}{-2.5mm}
    \item Số lớn là $24$, số bé là $6$:
    $$24 : 6 = 4\text{ (lần)}$$
    Ta nói: Số lớn gấp $4$ lần số bé.
    \item Số lớn là $35$, số bé là $7$: $35 : 7 = 5\text{ (lần)}$.
  \end{itemize}
\end{itemize}
}
\end{tcolorbox}
\end{minipage}
\hfill
\begin{minipage}[t]{0.49\textwidth}\vspace{0pt}
\begin{tcolorbox}[
  enhanced,
  colback=teal!4!white,
  colframe=teal!75!black,
  arc=2mm, boxrule=0.8pt,
  left=2mm, right=2mm, top=1.8mm, bottom=1.8mm,
  title={\footnotesize\bfseries\color{white}\faBalanceScale\ 2. Mối liên hệ hai chiều thuận nghịch}
]
{\footnotesize
\begin{itemize}\setlength{\itemsep}{1.5pt}\setlength{\leftskip}{-2.5mm}
  \item Nếu \textbf{số lớn gấp $n$ lần số bé}, thì:
  $$\implies \text{Số bé bằng } \mathbf{\dfrac{1}{n}} \text{ số lớn}$$
  \item \textbf{Ví dụ liên hệ:}
  \begin{itemize}\setlength{\itemsep}{0.5pt}\setlength{\leftskip}{-2.5mm}
    \item Mẹ $36$ tuổi, con $9$ tuổi.
    \item Tuổi mẹ gấp tuổi con: $36 : 9 = 4$ (lần).
    \item Tuổi con bằng $\mathbf{\dfrac{1}{4}}$ tuổi mẹ.
  \end{itemize}
  \item Đơn vị của câu trả lời là \textbf{"lần"}, không phải đơn vị đo đại lượng!
\end{itemize}
}
\end{tcolorbox}
\end{minipage}

\vspace{-2.5mm}
% --- KHỐI 3: SƠ ĐỒ TƯ DUY TIKZ ---
\noindent
\begin{tcolorbox}[
  enhanced,
  colback=orange!4!white,
  colframe=orange!80!black,
  arc=2mm, boxrule=0.8pt,
  left=2mm, right=2mm, top=1.5mm, bottom=1.5mm,
  title={\footnotesize\bfseries\color{white}\faSitemap\ SƠ ĐỒ TƯ DUY: CẶP QUAN HỆ SỐ LỚN -- SỐ BÉ}
]
\centering
\begin{tikzpicture}[
  arr/.style={-{Stealth[length=2mm]}, thick}
]
  \node[draw=blue!80!black, fill=blue!12, rounded corners=2pt, font=\scriptsize\bfseries, minimum width=3.2cm, minimum height=0.75cm, align=center] (b_lon) at (1.5, 0) {SỐ LỚN\\(Đoạn thẳng dài $a$)};
  \node[draw=teal!80!black, fill=teal!12, rounded corners=2pt, font=\scriptsize\bfseries, minimum width=3.2cm, minimum height=0.75cm, align=center] (b_be) at (8.5, 0) {SỐ BÉ\\(Đoạn thẳng ngắn $b$)};

  \draw[arr, purple!80!black] (b_lon.north east) to[bend left=25] node[midway, above, font=\scriptsize\bfseries] {Số lớn gấp mấy lần? $\implies a : b$ (lần)} (b_be.north west);
  \draw[arr, orange!80!black] (b_be.south west) to[bend left=25] node[midway, below, font=\scriptsize\bfseries] {Số bé bằng một phần mấy? $\implies \dfrac{1}{a : b}$} (b_lon.south east);
\end{tikzpicture}
\end{tcolorbox}

\vspace{-2.5mm}
% --- KHỐI 4: VÍ DỤ & MẸO TRÁNH BẪY ---
\noindent
\begin{minipage}[t]{0.49\textwidth}\vspace{0pt}
\begin{tcolorbox}[
  enhanced,
  colback=green!4!white,
  colframe=green!60!black,
  arc=2mm, boxrule=0.8pt,
  left=2mm, right=2mm, top=1.8mm, bottom=1.8mm,
  title={\footnotesize\bfseries\color{white}\faCheckCircle\ Ví dụ thực hành chuẩn mực}
]
{\footnotesize
\textbf{Bài toán mẫu:}\\
Sợi dây dài $32\text{ cm}$, sợi dây ngắn dài $8\text{ cm}$. Hỏi:\\
a) Sợi dây dài gấp mấy lần sợi dây ngắn?\\
b) Sợi dây ngắn bằng một phần mấy sợi dây dài?\\
\textbf{Bài giải:}\\
a) Sợi dây dài gấp sợi dây ngắn số lần là:
$$32 : 8 = \mathbf{4}\text{ (lần)}$$
b) Sợi dây ngắn bằng $\mathbf{\dfrac{1}{4}}$ sợi dây dài.
}
\end{tcolorbox}
\end{minipage}
\hfill
\begin{minipage}[t]{0.49\textwidth}\vspace{0pt}
\begin{tcolorbox}[
  enhanced,
  colback=red!4!white,
  colframe=red!70!black,
  arc=2mm, boxrule=0.8pt,
  left=2mm, right=2mm, top=1.8mm, bottom=1.8mm,
  title={\footnotesize\bfseries\color{white}\faExclamationTriangle\ Mẹo tránh bẫy trắc nghiệm}
]
{\footnotesize
\begin{itemize}\setlength{\itemsep}{1.5pt}\setlength{\leftskip}{-2.5mm}
  \item \textbf{Bẫy nhầm "gấp mấy lần" với "nhiều hơn mấy đơn vị":}\\
  -- Hỏi "gấp mấy lần" $\to$ làm tính \textbf{CHIA} ($32 : 8 = 4\text{ lần}$).\\
  -- Hỏi "nhiều hơn mấy cm" $\to$ làm tính \textbf{TRỪ} ($32 - 8 = 24\text{ cm}$).
  \item \textbf{Bẫy ghi nhầm đơn vị:}\\
  Đáp số phép chia tìm số lần có đơn vị là \textbf{"lần"}, không phải "cm" hay "kg" hay "tuổi"!
\end{itemize}
}
\end{tcolorbox}
\end{minipage}


\end{document}
